% Metodologia -- protocolo de validação manual (parte do Pedro).

\subsection{Validação manual}

Sorteamos 60 repositórios da amostra (\texttt{df.sample(n=60, random\_state=42)}, após
ordená-los pelo nome) e, de cada um, 5 \textit{releases} com \textit{release} anterior e
\texttt{compare} disponível; um repositório tinha só 4, e a amostra ficou com 299. Os três
autores rotularam de forma independente, sem consultar uns aos outros nem a saída da
heurística: (i) o tipo do projeto (biblioteca, aplicação, CLI ou outro); (ii) se as
\textit{releases} representam entregas reais ao usuário; e (iii) se cada \textit{release} é
corretiva, lendo as \textit{release notes} e os \textit{commits}. Medimos a concordância pelo
kappa de Fleiss, interpretado pela escala de Landis e Koch~\cite{landis1977}
(Tabela~\ref{tab:kappa}). O consenso é a maioria simples; empates 1-1-1 iriam para reunião,
mas não houve nenhum.

\begin{table}[ht]
\centering
\caption{Concordância entre os três avaliadores na amostra-ouro.}
\label{tab:kappa}
\small
\begin{tabular}{lcccc}
\hline
\textbf{Dimensão} & \textbf{Itens} & \textbf{Divergências} & \textbf{$\kappa$} & \textbf{Concordância} \\
\hline
Tipo do projeto             & 60  & 9  & 0,835 & Quase perfeita \\
Releases são entregas reais & 60  & 1  & 0,745 & Substancial \\
Release corretiva           & 299 & 26 & 0,876 & Quase perfeita \\
\hline
\end{tabular}
\end{table}

O menor kappa, o de entregas reais, vem do desbalanceamento: o consenso pôs 59 dos 60
repositórios em ``sim'', e com isso uma única divergência derruba o coeficiente, embora a
concordância bruta seja de 98,3\%. O tipo do projeto teve mais divergências, em projetos que
entregam ao mesmo tempo um servidor e um SDK e no alcance da categoria ``outro''. Em
\textit{release} corretiva, as 26 divergências vêm de \textit{releases} que misturam correções
e novidades. Como a heurística é avaliada contra esse consenso, o kappa de 0,876 serve de teto
prático para ela.
